\documentclass[../../DocumentoCompleto.tex]{subfiles}
\begin{document}
\chapter{Notas introductorias}
\section{Problema de investigación}

Las quemaduras de segundo grado por fricción son lesiones que implican la destrucción de la epidermis y la lámina basal \cite{Fotso2021}, debido principalmente a accidentes de tránsito o caídas \cite{Tracy2023}. Estas heridas son de importancia en Colombia, en donde tan solo en el mes de enero de 2025 se reportaron 953 accidentes de tránsito que tuvieron lesionados \cite{Agencia2025}, de los cuales aproximadamente el 12 \% presentaron quemaduras por fricción \cite{Trujillo2018}. A pesar de su frecuencia, esta problemática ha sido ignorada, ya que están asociadas a lesiones mecánicas \cite{Agrawal2008}, lo que provoca que las causas del problema no sean corregidas y genera consecuencias como un mayor riesgo de infecciones y complicaciones en la herida \cite{Miranda2020}. Esto lleva a un mayor tiempo de hospitalización de pacientes y al uso prolongado de antibióticos. Además, debido a la falta de herramientas eficaces para el tratamiento de heridas, es posible que se presente saturación en el sistema de salud junto con riesgos en la remisión de pacientes. 

La primera intervención para estas heridas en urgencias implica el desbridamiento, el cual es un procedimiento de limpieza. En el caso de quemaduras que abarquen un área menor al 20 \% del cuerpo, consta de un lavado con jabón y agua caliente, depilación del vello de la zona y cubrimiento con gasa vaselinada \cite{Zhang2022}. Para heridas mayores, es necesario realizar un tratamiento más complejo: el mismo protocolo de desbridamiento y remoción del tejido mediante procedimientos quirúrgicos \cite{Kim2024}. En estos casos graves, la etapa final de la curación se realiza haciendo uso de apósitos \cite{Miranda2020}, los cuales permiten tratar este tipo de quemaduras promoviendo la cicatrización, controlando el dolor y generando un ambiente antimicrobiano \cite{Miranda2020}.

Sin embargo, el uso de estos dispositivos se ha visto limitado debido a los altos costos y al bajo desarrollo de nuevas tecnologías transferibles al ámbito clínico \cite{Martinez2020}. En 2017, el costo del tratamiento para la hidratación de quemaduras intermedias era superior a 100 millones de pesos por paciente \cite{Archila2017}. Además, el número de artículos sobre apósitos es mayor que el número de patentes en este campo \cite{Martinez2020}, lo cual conlleva a una cantidad reducida de opciones comerciales y desmotiva la búsqueda de nuevas alternativas. Asimismo, se evidenció que, en un grupo de personal de urgencias, el 86.5 \% presentaba un desempeño de bajo nivel en el uso de tecnología para el tratamiento de heridas \cite{Galvis2018}, debido a deficiencias en la capacitación del personal médico de urgencias, lo que genera dependencia de especialistas como cirujanos plásticos.

Para solucionar estos problemas, es necesario el desarrollo de un apósito enfocado en quemaduras de segundo grado por fricción, que pueda ser utilizado en urgencias después del desbridamiento y que promueva la cicatrización.

\section{Pregunta de investigación}
¿Cómo se puede diseñar un sistema para favorecer la epitelización después del desbridamiento en el tratamiento de quemaduras de segundo grado por fricción ?
\section{Objetivo General}
Desarrollar un apósito que favorezca la epitelización después del desbridamiento en el tratamiento de quemaduras de segundo grado por fricción.
\section{Objetivos Específicos}

\begin{enumerate}
    \item Identificar las especificaciones de los apósitos que favorezcan la epitelización después del desbridamiento en el tratamiento de quemaduras de segundo grado por fricción, a partir de los requerimientos del cliente, normativos y técnicos.

    \item Fabricar un prototipo de apósito con geometría auxética que favorezca la epitelización después del desbridamiento en el tratamiento de quemaduras de segundo grado por fricción.

    \item Evaluar el comportamiento mecánico y biológico \textit{in vitro} de apósitos con geometría auxética que favorezcan la epitelización después del desbridamiento en el tratamiento de quemaduras de segundo grado por fricción.
\end{enumerate}

\section{Estado del arte}

Para llevar a cabo el estado del arte, se realizó una búsqueda en diferentes bases de datos de publicaciones científicas, tales como ScienceDirect, MDPI, PubMed, Google Scholar, Medline y Scopus, priorizando artículos relevantes publicados desde 2017. Se verificó que estos estuvieran indexados en revistas Q1 o Q2, utilizando Scimago como herramienta de validación para asegurar la calidad y relevancia de la información recopilada. Las palabras clave empleadas fueron ``Apósito'', ``Quemadura por fricción'', ``Cicatrización'' y ``Fármaco'', limitando la búsqueda a artículos que incluyeran el término ``Quemadura de tercer grado''.

En la actualidad, las soluciones para el tratamiento de quemaduras por fricción son los injertos de piel y los apósitos \cite{Suca2024}. En particular, los injertos autólogos se obtienen mediante un procedimiento quirúrgico donde un trozo de piel sana es ubicado sobre el tejido afectado \cite{Bogdanov2021}. Por otro lado, los apósitos cubren la herida mientras promueven un ambiente húmedo y bacteriostático que favorece la cicatrización, previene infecciones y evita procedimientos dolorosos como los injertos de piel \cite{Aderibigbe2018}. Sin embargo, ambos tratamientos presentan limitaciones debido a su baja capacidad de expansión, lo que impide que cubran áreas extensas \cite{Gupta2023}.

A pesar de que se ha evidenciado que las estructuras auxéticas han sido utilizadas para mitigar esta limitación en injertos de piel, hasta el momento no se han implementado este tipo de geometrías para la fabricación de apósitos para quemaduras \cite{Gupta2023}. En este sentido, las geometrías auxéticas, caracterizadas por un coeficiente de Poisson negativo, tienen la capacidad de expandirse en dirección perpendicular cuando se estiran axialmente, permitiendo cubrir heridas de mayor área \cite{Dou2023}. Entre estas, se destacan las estructuras auxéticas tipo ``I'', tipo ``Y'' y tipo panal de abeja, denominadas según su patrón estructural \cite{GuptaChanda2023}. Teniendo como referencia la geometría tradicional que alcanza una relación de mallado de 1.5:1 \cite{GuptaChanda2023}, en comparación, la estructura tipo ``I'' se caracteriza por lograr expansión uniforme alcanzando una relación de mallado superior a 4:1. Por otro lado, la geometría tipo ``Y'' consta de patrones conectados en forma de estrella que bajo esfuerzo presentan una expansión media, con una relación de mallado de 3:1. Por último, la geometría tipo panal de abeja tiene una expansión baja con una relación de mallado de aproximadamente 2.2:1 \cite{Gupta2023}. Pese al vacío en la literatura sobre la aplicación de estas estructuras en apósitos, avances en biomateriales han demostrado la posibilidad de integrar mallas auxéticas obtenidas mediante bioimpresión 3D de sustratos de hidrogel de colágeno con células \cite{Wei2025}. Este reciente hallazgo permite plantear el desarrollo de apósitos con estructura auxética basada en hidrogeles, que además de adaptarse a áreas de gran extensión, puedan liberar fármacos de forma controlada, constituyendo así una opción innovadora en el tratamiento de quemaduras.

Con respecto al material empleado para la fabricación de los apósitos para el tratamiento de quemaduras por fricción, se ha evidenciado una transición desde enfoques pasivos hacia el desarrollo de sistemas bioactivos avanzados. Materiales naturales como la membrana amniótica, la gelatina y la piel de Tilapia han sido ampliamente estudiados por su alto contenido en colágeno tipo I \cite{Gaviria2018,Arauz2022}, lo que favorece la migración celular, la deposición de matriz extracelular y la formación de tejido de granulación \cite{Wang2025}. Adicionalmente, la presencia de citoquinas y factores de crecimiento en estos biomateriales promueve una reepitelización acelerada y menos dolorosa, mejorando la experiencia clínica y reduciendo la formación de tejido cicatricial \cite{Singh2011}. Por otra parte, la integración de materiales sintéticos ha permitido el diseño de apósitos inteligentes con funciones específicas y arquitectura controlada. Ejemplos recientes incluyen hidrogeles termoformables de gelatina-alginato impresos en 3D, que ofrecen porosidad regulable, alta capacidad de absorción de exudado y buena adaptabilidad a superficies anatómicas complejas \cite{Li2022}. También se han reportado nanofibras electrohiladas compuestas por polímeros biocompatibles como quitosano y PVP, cargadas con antioxidantes como la florizina, que reducen el estrés oxidativo, aceleran la reepitelización y modulan la respuesta inflamatoria \cite{Zhang2022}. Sin embargo, estas no aportan efecto antimicrobiano ni garantizan buena estabilidad estructural y citocompatibilidad, como sí lo hacen los nanocompuestos híbridos de quitosano con nanopartículas de plata \cite{Islam2023}. Asimismo, estrategias de diseño biomimético han dado lugar a apósitos compuestos que combinan superficies externas hidrofóbicas con capas internas hidrofílicas, lo que imita la arquitectura de la piel y genera un microambiente favorable para la regeneración tisular en condiciones dinámicas \cite{Wang2023,Chen2024}.

Una estrategia viable para mejorar la interacción entre el apósito y el paciente ha sido la incorporación de diferentes aditivos, tales como antibióticos, proteinasas, analgésicos y factores de crecimiento \cite{Rezvani2019}. Teniendo en cuenta lo anterior, los antibióticos comúnmente utilizados son la gentamicina, minociclina, vancomicina y sulfadiazina de plata \cite{Freitez2024}. De estos, la vancomicina logra una mayor reducción de la carga bacteriana frente a \textit{Acinetobacter baumannii}, mientras que la gentamicina presenta mayor capacidad para favorecer la reepitelización \cite{Nuutila2019}. Asimismo, las proteinasas, que son enzimas proteolíticas que degradan células dañadas y necróticas mediante la hidrólisis de enlaces peptídicos, generan una reducción del color y del edema \cite{DaSilva2017}. Dentro de ellas, se destaca la colagenasa debido a su capacidad para hidrolizar los componentes del tejido necrótico \cite{Patry2017}. En cuanto al control del dolor, se ha incorporado morfina y clorhidrato de lidocaína en la formulación de los apósitos, con el fin de proporcionar analgesia localizada y mejorar la comodidad del paciente durante el tratamiento \cite{Kowalski2024}. Por otro lado, los factores de crecimiento han demostrado acelerar el proceso de granulación y promover la regeneración tisular mediante la estimulación de la angiogénesis, especialmente cuando se emplean en conjunto factores como el transformante beta (TGF-$\beta$), el endotelial vascular (VEGF), los derivados de plaquetas (PDGF) y el epidérmico (EGF) \cite{Heydari2024}.

A partir de la revisión realizada, se evidenció la ausencia de estudios que evalúen el impacto del uso de geometrías auxéticas en apósitos para quemaduras, a pesar de los beneficios demostrados en injertos. Esta brecha permite plantear el desarrollo de un apósito con geometría auxética, capaz de adaptarse a zonas de pliegues y articulaciones, incorporando aditivos bioactivos que favorezcan la reepitelización y mejoren las condiciones de cicatrización.


\section{Marco Teórico}
  \subsection{Quemaduras de fricción}
 Una quemadura por fricción se produce cuando la piel entra en contacto repetido o forzado con una superficie, lo que ocasiona un daño que combina la abrasión mecánica con el calor generado por el roce \cite{Zuluaga2023}. Esta interacción puede originar desde lesiones superficiales, limitadas a la epidermis, hasta comprometer la dermis o capas más profundas si la fuerza y la duración del contacto son mayores \cite{Tracy2023}. Se trata de un tipo de lesión que aparece con frecuencia en accidentes de tránsito, especialmente en motociclistas y ciclistas que sufren deslizamientos sobre el pavimento \cite{Perkins2025}, aunque también se observa en deportes de contacto o caídas sobre suelos ásperos \cite{MacFarlane2021}. Las regiones más afectadas suelen ser aquellas expuestas y prominentes, como rodillas y codos, que tienden a absorber el impacto directo \cite{Perkins2024}.

  \subsection{Apósitos en el tratamiento de quemaduras}
El uso de apósitos en el manejo de las quemaduras constituye un recurso fundamental dentro del proceso de recuperación. Estos dispositivos intervienen directamente en la protección del tejido dañado y en la creación de un ambiente favorable para la cicatrización \cite{Noor2022}. Al mantener un grado de humedad adecuado y disminuir el riesgo de infecciones, los apósitos facilitan el proceso de epitelización y contribuyen a una regeneración más eficiente de la piel \cite{Opriessnig2023}.  

\par  
Entre las distintas opciones disponibles, los hidrogeles destacan por sus propiedades. Debido a su elevado contenido de agua, refrescan de manera inmediata la superficie de la quemadura, lo que proporciona alivio al paciente y previene la resequedad \cite{Shu2021}. Además, generan un ambiente húmedo y estable que favorece la reparación tisular y estimula la formación de nuevo epitelio, razón por la cual se consideran una alternativa de elección en quemaduras superficiales y de espesor parcial \cite{Celik2024}.


  \subsection{Coeficiente de Poisson}
El coeficiente de Poisson indica la relación entre la deformación lateral y la deformación longitudinal que experimenta un material cuando se le aplica un esfuerzo uniaxial \cite{Farzaneh2022}. En este sentido, refleja cómo varía la dimensión transversal de un cuerpo respecto a su cambio en la dimensión longitudinal y se expresa como la razón negativa entre ambas deformaciones, tal como se muestra en la ecuación~(\ref{poisson}) \cite{Chang2022}.


\begin{equation}
\nu = - \frac{\varepsilon_{t}}{\varepsilon_{l}}
\label{poisson}
\end{equation}

  \subsection{Geometría auxética}
 Las estructuras que presentan un coeficiente de Poisson negativo se conocen como geometrías auxéticas \cite{Ebert2024}. Su comportamiento mecánico está determinado principalmente por la disposición de su arquitectura interna y no únicamente por la composición del material \cite{Yolcu2022}. Por lo tanto, cuando se someten a una fuerza de tracción unidireccional, estas estructuras se expanden en dirección transversal, un comportamiento que contrasta con la respuesta habitual de las geometrías convencionales \cite{Winczewski2022}.
  
  \subsection{Agarosa}
  La agarosa es un polisacárido lineal de origen natural, obtenido a partir de algas rojas como \textit{Gelidium} y \textit{Gracilaria} \cite{Mantri2021}. Su estructura está compuesta por unidades repetitivas de agarobiosa \cite{Roux2023}, un disacárido integrado por D-galactosa y 3,6-anhidro-L-galactosa unidas mediante enlaces glicosídicos \cite{Shen2024}. Una de sus propiedades más importantes es la capacidad de formar geles termo-reversibles en soluciones acuosas, generando una red porosa y estable. Estas características han hecho de la agarosa un material ampliamente utilizado en biología molecular, en la fabricación de hidrogeles para aplicaciones biomédicas y en sistemas de liberación controlada de fármacos \cite{Jiang2023}.
  
  \subsection{Ácido hialurónico}
El ácido hialurónico (AH) es un polisacárido lineal de la familia de los glucosaminoglucanos (GAGs) \cite{Yasin2022}, compuesto por unidades repetitivas de ácido D-glucurónico y N-acetil-D-glucosamina, unidas mediante enlaces $\beta$(1$\to$3) y $\beta$(1$\to$4) \cite{Lee2021}. Se encuentra de manera natural en la matriz extracelular de los tejidos animales, donde cumple funciones esenciales en la hidratación, la lubricación y la organización estructural \cite{Ding2022}.

  \subsection{Sulfadiazina de plata}
La sulfadiazina de plata (SSD) es un antibiótico el cual combina sulfonamida (sulfadiazina) con iones de plata \cite{Nunes2023}, unión que le confiere una potente acción antimicrobiana. La sulfadiazina actúa inhibiendo la síntesis de ácido fólico bacteriano, lo que limita la multiplicación de los microorganismos \cite{Fatima2022}, mientras que la plata ejerce un efecto bactericida de amplio espectro al alterar las membranas celulares y desestabilizar el material genético de los patógenos. Esta acción conjunta hace de la sulfadiazina de plata un fármaco eficaz en la prevención y control de infecciones cutáneas extensas \cite{DeFrancesco2022}.

\end{document}
